\section{Extended Abstract}\label{sec:extended_abstract}

This paper advocates the use of mobile phones or tablets for blockchain mining.
Namely, they form a large and cheap hardware base, which is
interesting, in particular, to start up a new blockchain, when it is
hard to convince people to install complex software
on a continuously connected computer. This paper acknowledges that mobile phones
and tablets have hardware limitations that make most mining algorithms impractical,
with the notable exception of proof of space.

Namely, proof of space~\cite{AbusalahACKPR17,AtenieseBFG14,CohenP19,DziembowskiFKP15,ParkKFGAP18,RenD16,Spoto25}
is a mining algorithm where disk (currently SSD) space
is used instead of CPU power.
With proof of space, blocks carry evidence
of the allocation of a large data structure on disk for their creation.
The larger that structure, the higher the likelihood that
blocks get included in the stable chain. Proof of space is decentralized like
Bitcoin's proof of work and ecological like proof of stake.
More specifically, this paper builds on the proof of space algorithm formalized in~\cite{Spoto25},
where a preliminary, expensive
\emph{plotting phase} allocates large data on disk,
which is later used in a
cheap, inexpensive \emph{mining phase} to quickly answer
block creation challenges received from blockchain peers.

Mobile mining with a mobile phone or tablet is appealing because:
%
\begin{itemize}
\item it leverages a huge existing hardware base of devices, that otherwise remain idle
  for most of the time;
\item it introduces mining to a large portion of humanity, who owns a
  phone and a limited internet connection, but not a powerful, well-connected computer;
\item it can be easily updated, thanks to the automatic update system of mobile devices;
\item it simplifies the start-up of new blockchains, since it is easier to convince people
  to install an app than to run and maintain a software on an always connected computer.
\end{itemize}
%
Despite that, mobile mining has not been used yet because of obvious limitations
of mobile devices \wrt{} CPU power, battery capacity and network bandwidth, in particular
in least favorite countries.
This paper shows that
proof of space emerges as a niche solution here, whose mining
is largely insensitive to these limitations.

This paper makes the following contributions:
%
\begin{enumerate}
\item it presents mobile mining, its power and limitations;
\item it shows how proof of space emerges as a niche solution here, insensitive
  to such limitations;
\item it presents Mokaminter, our actual Android app for mobile mining with proof of space;
\item it reports experiments with Mokaminter, showing its extremely low
  CPU use, memory use, battery consumption and amount of exchanged data.
\end{enumerate}

Mokaminter runs only a miner in the phone, not a full blockchain
peer, which would be prohibitive for the size of the blockchain,
that of its state database, the amount of data exchanged and the computational
cost. The blockchain peer is an external computer, to which the phone connects.

For the best of our knowledge, Mokaminter is the only existing software able
to mine with a mobile phone or tablet. Namely,
a simple search on Google Play reports a few Android apps for
\emph{crypto mining}. Despite their name, these apps are not miners, but either:
mobile dashboards to cloud pool miners (\ie{}, mining occurs in the cloud and
the app is just monitoring)
or dubious games that promise Bitcoins in exchange for solving tasks
or watching ads.

Mokaminter is an open-source application~\cite{mokaminter-source-code}
distributed on Google Play~\cite{mokaminter-google-play}
under the terms of the Apache 2.0 License. It mines for all
blockchains based on the Mokamint generic tool for building
blockchains based on proof of space~\cite{mokamint-source-code}. Currently,
the Hotmoka blockchain~\cite{hotmoka} is
built over Mokamint, therefore Mokaminter can be used for mining for Hotmoka.
